%! Unit 3, Lessons 3.1–3.5 — SAMPLE Test (Answer Key)
% The practice form of the Unit 3 test. Item for item it is the parallel form of ../actual_test/:
% same two sections, the same ten multiple-choice targets in the same order, the same two
% free-response sets with the same parts and the same points (MC 10, FRQ 1 = 4, FRQ 2 = 6,
% total 20) — every context different, so working this one teaches the moves, not the answers.
% Merged into the unit student packet via ../../sample_test/ (see shared/tests.mk).
% NAMESTRIP: no name line — mirrors the blank (a graded page; the name is on the cover sheet only).
% NO teachernote here — scoring guidance lives in ../../rubrics/, never in a key.
% GENERATED FROM the blank: the correct MC options wrapped in \ans{}, each
% \gradedhide{...\writelines{n}} turned into the box plus n \ansline{}.
% NO SKETCH ITEMS: the one display (the reshuffle dotplot) is pre-drawn in TikZ and read.
%
% BLUEPRINT (two MC per lesson, in lesson order; the FRQs carry the written moves)
%  MC1  3.1 observational vs experiment — careful measurement is not a treatment (the crux)
%  MC2  3.1 generalize only to the population sampled
%  MC3  3.2 name the sampling method (stratified here; cluster on the actual test)
%  MC4  3.2 equal chance for each individual is not an SRS (one whole grade at random)
%  MC5  3.3 a random sample can still be biased (undercoverage here; response bias on the actual)
%  MC6  3.3 a larger sample does not fix bias (voluntary response)
%  MC7  3.4 why random assignment: balance confounding variables
%  MC8  3.4 placebo and blinding
%  MC9  3.4 matched pairs — chance still decides inside every pair
%  MC10 3.5 scope of inference: random assignment, no random selection
%  FRQ1 3.2/3.3 SRS by random number generator; cluster vs stratified; nonresponse + direction
%  FRQ2 3.4/3.5 units/treatments/response; random assignment; blocking; reading a reshuffle
%       dotplot; cause and to whom.
%
% THE DATA
%  FRQ2(d) 100 reshuffles, difference in mean words (music - silence), step 0.3:
%    -2.1:1 -1.8:1 -1.5:3 -1.2:6 -0.9:9 -0.6:12 -0.3:14 0:14 0.3:13 0.6:11 0.9:8 1.2:4
%    1.5:2 1.8:1 2.1:1  (sum 100).  Observed 1.8: 2 of 100 at or above -> significant.
%
% PAGE LOCKSTEP with ../../test_keys/practice_test_key/: every \writelines{n} here is answered
% by exactly n \ansline{} there, each terse enough not to wrap. Nothing else differs except the
% preamble, the header, and the \ans{} marks on the correct MC options.
\documentclass[12pt]{article}
\usepackage{apstats-article}
\usepackage{apstats-key}   % pulls in -boxes; do NOT also load -boxes
\usepackage{grading}
\gradingassignment{apstats-u03-practice-test}
\gradingcourse{advanced_statistics}

\begin{document}

\pageheader{Unit 3 \quad Lessons 3.1--3.5}{Sample Test --- Answer Key}
{\small\itshape Write each final answer in its box. Only what is inside a box is graded.}\par

{\small\hfill Score:\ \blank{1.4cm}\ / 20}

\vspace{0.06in}

\boxguard[8]
\begin{remindbox}
\small
This is the \textbf{practice} test --- not scored. Work it in one period, notes closed, then check
it against the key. It matches the real test item for item; every answer must be in context.
\end{remindbox}

\vspace{0.10in}

\boxguard[16]
\begin{headlinebox}{goldbg}
\small
\textbf{Section I --- Multiple Choice} \hfill \textbf{10 points, 1 each}\\
Write the letter of the single best answer in its Letter box. Read all five options --- several
name a real idea from the unit, but only one fits \emph{this} study.
\end{headlinebox}

\vspace{0.10in}

\begin{enumerate}[leftmargin=1.9em, itemsep=0.13in]

  \boxguard[11]
  \item Researchers gave 300 adults a wrist tracker that recorded every minute of their sleep for
        six weeks, and each adult logged every cup of coffee in an app. Adults who drank more
        coffee slept less, on average. Which statement is correct?
        \begin{enumerate}[label=\textbf{(\Alph*)}, leftmargin=2.2em, itemsep=2pt]
          \item It is an experiment, because the tracker measured sleep precisely.
          \item It is an experiment, because the researchers gave every adult a tracker.
          \item \ans{It is an observational study, because the researchers did not assign how much coffee anyone drank.}
          \item It is an observational study, because it involved fewer than 1{,}000 adults.
          \item It is a census, because every minute of sleep was recorded.\hfill\gradedright{\textit{Letter:}~\gradedblank[id=t1, type=choice, points=1, answer=C]{0.45in}{0.2in}}\rule[-0.2in]{0pt}{0.48in}
        \end{enumerate}

  \boxguard[11]
  \item A state parks agency randomly selected 400 of the 25{,}000 people who bought an annual
        pass to Lake Harmon State Park last year and asked whether they support a new parking
        fee. To which group can the result be generalized?
        \begin{enumerate}[label=\textbf{(\Alph*)}, leftmargin=2.2em, itemsep=2pt]
          \item All adults who live in the state.
          \item Everyone who visits any state park.
          \item Only the 400 pass holders who were asked.
          \item \ans{All people who bought an annual pass to Lake Harmon State Park last year.}
          \item No one, because a sample is not a census.\hfill\gradedright{\textit{Letter:}~\gradedblank[id=t2, type=choice, points=1, answer=D]{0.45in}{0.2in}}\rule[-0.2in]{0pt}{0.48in}
        \end{enumerate}

  \boxguard[12]
  \item A bank sorts its 9{,}000 customers by account type --- checking only, savings only, or
        both --- and takes a simple random sample of 100 customers from each type to estimate
        how satisfied its customers are. Which best describes this sample?
        \begin{enumerate}[label=\textbf{(\Alph*)}, leftmargin=2.2em, itemsep=2pt]
          \item \ans{A stratified random sample, because an SRS is taken from every account type.}
          \item A cluster sample, because the customers are split into groups.
          \item A simple random sample, because every customer had a chance to be chosen.
          \item A cluster sample, because every account type is included.
          \item A convenience sample, because the bank used its own customer list.\hfill\gradedright{\textit{Letter:}~\gradedblank[id=t3, type=choice, points=1, answer=A]{0.45in}{0.2in}}\rule[-0.2in]{0pt}{0.48in}
        \end{enumerate}

  \boxguard[13]
  \item A high school has exactly 400 students in each of grades 9, 10, 11, and 12. The
        principal picks one grade at random, then takes a simple random sample of 100 students
        from that grade. Every student has the same $\tfrac{1}{4}\cdot\tfrac{100}{400} =
        \tfrac{1}{16}$ chance of being chosen. Is this a simple random sample of 100 students
        from the school?
        \begin{enumerate}[label=\textbf{(\Alph*)}, leftmargin=2.2em, itemsep=2pt]
          \item Yes, because every student has the same 1-in-16 chance of being chosen.
          \item Yes, because chance is used at both steps.
          \item No, because the four grades are different sizes.
          \item No, because a simple random sample must be chosen with replacement.
          \item \ans{No: a sample with students from two grades can never be chosen.}\hfill\gradedright{\textit{Letter:}~\gradedblank[id=t4, type=choice, points=1, answer=E]{0.45in}{0.2in}}\rule[-0.2in]{0pt}{0.48in}
        \end{enumerate}

  \boxguard[14]
  \item A cable company wants to estimate the proportion of adults in its city who stream TV
        over the internet. It chooses 800 numbers at random from the city's \textbf{landline}
        telephone directory and calls them. Which problem is built into the way this sample was
        chosen?
        \begin{enumerate}[label=\textbf{(\Alph*)}, leftmargin=2.2em, itemsep=2pt]
          \item Voluntary response bias: the adults decide whether to be in the sample.
          \item \ans{Undercoverage bias: adults with no landline --- likely heavier streamers ---
                cannot be chosen, so the estimate is probably too low.}
          \item Nonresponse bias: some of the adults who are called will not answer.
          \item No bias, because the 800 numbers were chosen at random.
          \item Response bias: adults may misremember their streaming.\hfill\gradedright{\textit{Letter:}~\gradedblank[id=t5, type=choice, points=1, answer=B]{0.45in}{0.2in}}\rule[-0.2in]{0pt}{0.48in}
        \end{enumerate}

  \boxguard[14]
  \item A gaming website invited visitors to vote on whether a new game is ``too hard'':
        52{,}000 voted, and 74\% said yes. A researcher asked a random sample of 500 of the
        game's registered players the same question, and 38\% said yes. Which statement is
        best?
        \begin{enumerate}[label=\textbf{(\Alph*)}, leftmargin=2.2em, itemsep=2pt]
          \item The website poll is more trustworthy, because its sample is over 100 times larger.
          \item The two results differ only because of sampling variability.
          \item \ans{The random sample is more trustworthy: the website poll has voluntary response
                bias, and a larger sample does not fix bias.}
          \item Averaging the two, about 56\% of players think the game is too hard.
          \item Neither result is useful, because neither one is a census.\hfill\gradedright{\textit{Letter:}~\gradedblank[id=t6, type=choice, points=1, answer=C]{0.45in}{0.2in}}\rule[-0.2in]{0pt}{0.48in}
        \end{enumerate}

  \boxguard[14]
  \item A gym wants to know whether a 10-minute stretching routine before workouts reduces
        injuries. It will randomly assign 120 member volunteers to stretch or not to stretch
        for six months. What is the main purpose of the random assignment?
        \begin{enumerate}[label=\textbf{(\Alph*)}, leftmargin=2.2em, itemsep=2pt]
          \item So the results can be generalized to all gym members in the country.
          \item To make sure each group has exactly 60 members.
          \item To remove the need for a group that does not stretch.
          \item \ans{To make the two groups roughly balanced on variables such as age and fitness,
                so a difference in injuries can be traced to the stretching.}
          \item To guarantee that the stretching group has fewer injuries.\hfill\gradedright{\textit{Letter:}~\gradedblank[id=t7, type=choice, points=1, answer=D]{0.45in}{0.2in}}\rule[-0.2in]{0pt}{0.48in}
        \end{enumerate}

  \boxguard[15]
  \item To test whether a new nasal spray relieves allergy symptoms, 80 patients are randomly
        assigned to the new spray or to a salt-water spray that looks and smells the same.
        Neither the patients nor the doctors who rate their symptoms know who got which spray.
        Which statement is correct?
        \begin{enumerate}[label=\textbf{(\Alph*)}, leftmargin=2.2em, itemsep=2pt]
          \item \ans{The salt-water spray is a placebo, and the experiment is double-blind.}
          \item The salt-water spray is a placebo, and the experiment is single-blind.
          \item The salt-water group is not needed, because the doctors rate symptoms carefully.
          \item The experiment is double-blind because it compares two treatments.
          \item The salt-water spray is a confounding variable.\hfill\gradedright{\textit{Letter:}~\gradedblank[id=t8, type=choice, points=1, answer=A]{0.45in}{0.2in}}\rule[-0.2in]{0pt}{0.48in}
        \end{enumerate}

  \boxguard[15]
  \item Each of 30 volunteers will spread sunscreen A on one forearm and sunscreen B on the
        other, spend two hours in the sun, and have the redness of each forearm rated. Most
        people use their right arm more, and it may get more sun. Which plan is best?
        \begin{enumerate}[label=\textbf{(\Alph*)}, leftmargin=2.2em, itemsep=2pt]
          \item Put sunscreen A on every volunteer's right forearm and B on the left.
          \item Let each volunteer choose which forearm gets sunscreen A.
          \item Each volunteer wears both sunscreens, so no randomization is needed.
          \item Randomly assign 15 volunteers to wear A on both arms and 15 to wear B on both.
          \item \ans{For each volunteer, a coin flip decides which forearm gets A.}\hfill\gradedright{\textit{Letter:}~\gradedblank[id=t9, type=choice, points=1, answer=E]{0.45in}{0.2in}}\rule[-0.2in]{0pt}{0.48in}
        \end{enumerate}

  \boxguard[15]
  \item A researcher recruits 150 volunteers through a newspaper ad and randomly assigns 75 to a
        daily 10-minute meditation and 75 to a daily 10-minute reading break. After eight weeks
        the meditation group reports lower stress, and the difference is statistically
        significant. Which conclusion is most appropriate?
        \begin{enumerate}[label=\textbf{(\Alph*)}, leftmargin=2.2em, itemsep=2pt]
          \item Meditation causes lower stress for all adults.
          \item \ans{Meditation causes lower stress for volunteers like these, but the result cannot
                be generalized to a larger population.}
          \item Among all newspaper readers, meditators tend to report lower stress, but cause
                cannot be concluded.
          \item No conclusion is possible, because the subjects were volunteers.
          \item Meditation causes lower stress for all readers of the newspaper.\hfill\gradedright{\textit{Letter:}~\gradedblank[id=t10, type=choice, points=1, answer=B]{0.45in}{0.2in}}\rule[-0.2in]{0pt}{0.48in}
        \end{enumerate}

\end{enumerate}

\vspace{0.14in}

% No \newpage here: the guard keeps the Section II strip off a page edge.
\boxguard[30]
\begin{headlinebox}{goldbg}
\small
\textbf{Section II --- Free Response} \hfill \textbf{10 points}\\
Show your reasoning. Answer \emph{in context} --- name the people or things studied and the
variable measured.
\end{headlinebox}

\vspace{0.10in}

\begin{enumerate}[leftmargin=1.9em, itemsep=0.16in, resume]

  \boxguard[26]
  \item \textbf{(4 points)} A city library system has 12{,}000 cardholders, registered at its 8
        branches. It wants to estimate the proportion of cardholders who would use a new
        library phone app.

        \vspace{6pt}
        \textbf{(a)} Describe how to use a random number generator to choose a simple random
        sample of 200 cardholders.\\[4pt]
        \gradedhere{\gradedblank[id=t11a, type=rubric, points=1, rubric=../rubrics/apstats-u03-practice-test-t11a.yaml]{\linewidth+6pt}{\gradedlinesheight{3}}}\ansline{Label the cardholders 1 to 12{,}000. Use the generator to produce random}\\
        \ansline{integers from 1 to 12{,}000, skipping any repeat, until 200 different labels}\\
        \ansline{are chosen. Survey the 200 cardholders with those labels.}\\
        \ansline{}\\
        \ansline{}\\
        \ansline{}\\
        \ansline{}\\
        \ansline{}

        \boxguard[16]
        \vspace{6pt}
        \textbf{(b)} \emph{Plan A:} randomly choose 2 of the 8 branches and survey every
        cardholder registered there. \emph{Plan B:} split the cardholders into three age
        groups --- under 30, 30 to 59, and 60 or older --- and take a simple random sample from
        each. Name the sampling method each plan uses, and explain why Plan B should give a
        more precise estimate.\\[4pt]
        \gradedhere{\gradedblank[id=t11b, type=rubric, points=1, rubric=../rubrics/apstats-u03-practice-test-t11b.yaml]{\linewidth+6pt}{\gradedlinesheight{3}}}\ansline{Plan A is a cluster sample; Plan B is a stratified random sample.}\\
        \ansline{App use is likely related to age, so each age group is alike inside. An SRS}\\
        \ansline{from every group guarantees each is represented, so the estimate varies less.}\\
        \ansline{Two random branches could happen to hold mostly older (or younger) users.}\\
        \ansline{}\\
        \ansline{}\\
        \ansline{}\\
        \ansline{}

        \boxguard[16]
        \vspace{6pt}
        \textbf{(c)} The library e-mails the survey to the 200 cardholders in its simple random
        sample, and only 74 reply. Name the bias this creates. Explain who is likely to be
        over-represented, and whether the estimate is likely too high or too low.\\[4pt]
        \gradedhere{\gradedblank[id=t11c, type=rubric, points=2, rubric=../rubrics/apstats-u03-practice-test-t11c.yaml]{\linewidth+6pt}{\gradedlinesheight{3}}}\ansline{Nonresponse bias: 126 of the 200 chosen cardholders did not answer.}\\
        \ansline{Cardholders who answer an e-mail survey are likely more comfortable with}\\
        \ansline{e-mail and phones, so likely app users are over-represented. The estimate}\\
        \ansline{of the proportion who would use the app is likely too high.}\\
        \ansline{}\\
        \ansline{}\\
        \ansline{}\\
        \ansline{}

  \boxguard[24]
  \item \textbf{(6 points)} A teacher wants to know whether listening to classical music while
        studying improves scores on a 20-word vocabulary quiz. Forty student volunteers will be
        randomly assigned, 20 to study the word list for 15 minutes with classical music
        playing and 20 to study it in silence. Then every student takes the quiz.

        \vspace{6pt}
        \textbf{(a)} Identify the experimental units, the treatments, and the response
        variable.\\[4pt]
        \gradedhere{\gradedblank[id=t12a, type=rubric, points=1, rubric=../rubrics/apstats-u03-practice-test-t12a.yaml]{\linewidth+6pt}{\gradedlinesheight{3}}}\ansline{Units: the 40 student volunteers.}\\
        \ansline{Treatments: studying with classical music; studying in silence.}\\
        \ansline{Response: the score on the 20-word vocabulary quiz.}\\
        \ansline{}\\
        \ansline{}\\
        \ansline{}\\
        \ansline{}\\
        \ansline{}

        \boxguard[14]
        \vspace{6pt}
        \textbf{(b)} Describe how the teacher could use 40 identical slips of paper to carry out
        the random assignment.\\[4pt]
        \gradedhere{\gradedblank[id=t12b, type=rubric, points=1, rubric=../rubrics/apstats-u03-practice-test-t12b.yaml]{\linewidth+6pt}{\gradedlinesheight{3}}}\ansline{Write each student's name on one of the 40 identical slips, mix them well}\\
        \ansline{in a hat, and draw 20 without replacement. Those 20 study with music; the}\\
        \ansline{other 20 study in silence.}\\
        \ansline{}\\
        \ansline{}\\
        \ansline{}\\
        \ansline{}\\
        \ansline{}

        \boxguard[15]
        \vspace{6pt}
        \textbf{(c)} The teacher believes a student's reading level (advanced or regular)
        affects quiz scores. Explain how to block on reading level, and what chance must still
        decide inside each block.\\[4pt]
        \gradedhere{\gradedblank[id=t12c, type=rubric, points=1, rubric=../rubrics/apstats-u03-practice-test-t12c.yaml]{\linewidth+6pt}{\gradedlinesheight{3}}}\ansline{Split the students into two blocks: advanced readers and regular readers.}\\
        \ansline{Inside each block, chance (slips in a hat) assigns half the students to}\\
        \ansline{music and half to silence. Compare music to silence within each block.}\\
        \ansline{}\\
        \ansline{}\\
        \ansline{}\\
        \ansline{}\\
        \ansline{}

        \boxguard[30]
        \vspace{6pt}
        \textbf{(d)} The music group's mean score was \textbf{1.8 words higher} than the silent
        group's. A computer reshuffled the 40 scores into two new groups of 20, 100 times, as
        if the music made no difference. How many reshuffles gave a difference of 1.8 words or
        more? Is the observed difference statistically significant? Explain.
        \begin{center}
        \begin{tikzpicture}[x=0.75cm, y=0.1cm]
          \draw[thick] (-7.5,0) -- (7.5,0);
          \foreach \x/\l in {-7/-2.1,-6/-1.8,-5/-1.5,-4/-1.2,-3/-0.9,-2/-0.6,-1/-0.3,0/0,
                             1/0.3,2/0.6,3/0.9,4/1.2,5/1.5,6/1.8,7/2.1} {
            \draw (\x,0) -- (\x,-1.4); \node[below, font=\tiny] at (\x,-1.2) {$\l$};}
          \foreach \x/\n in {-7/1,-6/1,-5/3,-4/6,-3/9,-2/12,-1/14,0/14,1/13,2/11,3/8,4/4,
                             5/2,6/1,7/1}
            \foreach \k in {1,...,\n} \fill[navy] (\x,\k) circle (1.3pt);
          \draw[redacc, dashed, thick] (6,3) -- (6,14);
          \node[font=\tiny\bfseries, redacc, above] at (6,14) {observed 1.8};
          \node[font=\tiny] at (0,-7) {Simulated difference in mean words correct, music $-$ silence};
          \node[font=\scriptsize\bfseries] at (0,19) {100 random reshuffles};
        \end{tikzpicture}
        \end{center}
        \gradedhere{\gradedblank[id=t12d, type=rubric, points=2, rubric=../rubrics/apstats-u03-practice-test-t12d.yaml]{\linewidth+6pt}{\gradedlinesheight{3}}}\ansline{2 of the 100 reshuffles gave a difference of 1.8 words or more.}\\
        \ansline{Yes, it is significant: a difference this large almost never happens by the}\\
        \ansline{luck of the random assignment alone, so it is not easily explained by chance.}\\
        \ansline{}\\
        \ansline{}\\
        \ansline{}\\
        \ansline{}\\
        \ansline{}

        \boxguard[15]
        \vspace{6pt}
        \textbf{(e)} Can the teacher conclude that the classical music \emph{caused} the higher
        scores? To whom can the result be generalized? Explain.\\[4pt]
        \gradedhere{\gradedblank[id=t12e, type=rubric, points=1, rubric=../rubrics/apstats-u03-practice-test-t12e.yaml]{\linewidth+6pt}{\gradedlinesheight{3}}}\ansline{Yes: the treatments were randomly assigned and the difference is significant,}\\
        \ansline{so music caused the higher scores. But the students were volunteers, not a}\\
        \ansline{random sample, so the result applies only to students like these 40.}\\
        \ansline{}\\
        \ansline{}\\
        \ansline{}\\
        \ansline{}\\
        \ansline{}

\end{enumerate}

\end{document}
